% rhythm: beamer frames under moloch, top-aligned: frame title to the first body line, paragraph to paragraph, list items, block title to body, two columns' first lines, and where a centred frame's one line and the frame number stand
% boundary: frametitle-body Birch span
% boundary: par-par Cedar span
% boundary: par-list Dogwood span
% boundary: item-item Elm span
% boundary: list-par Hazel span
% boundary: frametitle-block Larch span
% boundary: blocktitle-body Maple span
% boundary: block-par Oak span
% boundary: column-top Rowan top
% boundary: column-top Spruce top
% boundary: centred-top Walnut top
% boundary: page-footline =4 top
\documentclass{beamer}
\usetheme{moloch}
\usepackage{fontspec}
\setsansfont{OpenSans-Regular.ttf}[Path=../corpus/fonts/, BoldFont=OpenSans-Bold.ttf, ItalicFont=OpenSans-Italic.ttf, BoldItalicFont=OpenSans-BoldItalic.ttf]
\setmainfont{OpenSans-Regular.ttf}[Path=../corpus/fonts/, BoldFont=OpenSans-Bold.ttf, ItalicFont=OpenSans-Italic.ttf, BoldItalicFont=OpenSans-BoldItalic.ttf]
\begin{document}
\begin{frame}[t]{Alder title}
Birch is the frame's first paragraph.

Cedar is its second paragraph.
\begin{itemize}
\item Dogwood is the first item.
\item Elm is the second item.
\end{itemize}
Hazel is the paragraph after the list.
\end{frame}
\begin{frame}[t]{Juniper title}
\begin{block}{Larch block}
Maple is the block's body.
\end{block}
Oak is the paragraph after the block.
\end{frame}
\begin{frame}[t]{Pine title}
\begin{columns}[T]
\begin{column}{0.45\textwidth}
Rowan is the left column.
\end{column}
\begin{column}{0.45\textwidth}
Spruce is the right column.
\end{column}
\end{columns}
\end{frame}
\begin{frame}{Teak title}
Walnut is the one line of a centred frame.
\end{frame}
\end{document}
